% !PRINTED 747
\weberpage{747}
Um diesen wichtigen Satz nachzuweisen, erinnern wir uns daran, dass jede algebraische Zahl $\Theta$ die Wurzel einer und nur einer irreduciblen Gleichung
\weq{1}{f(\Theta) = a_0 \Theta^n + a_1 \Theta^{n-1} + \cdots + a_n = 0}
ist, in der $a_0, a_1, \ldots, a_n$ ganze rationale Zahlen ohne gemeinsamen Theiler sind, und $a_0$ von Null verschieden und positiv ist. Der Grad $n$ der Gleichung (1) ist eine positive ganze Zahl, also mindestens $=1$. Die Zahlen $a_1, \ldots, a_{n-1}$ können zum Theil Null sein.

Wir wollen nun die Vorzeichen $\pm$ so bestimmen, dass $\pm a_1, \pm a_2, \pm a_3, \ldots, \pm a_n$ nicht negativ sind, und nennen die Summe
\weq{2}{N = (n-1) + a_0 \pm a_1 \pm a_2 \cdots \pm a_n}
die \gesperrt{Höhe} der algebraischen Zahl $\Theta$. Die Höhe ist dann immer eine \gesperrt{positive} ganze Zahl.

Nun ist leicht zu sehen, dass es für einen gegebenen Werth der Höhe $N$ immer nur eine endliche Anzahl von algebraischen Zahlen geben kann. Denn zunächst kann $n$ nach (2) niemals grösser als $N$ sein, und zu jedem gegebenen $N$ und $n$ können die Zahlen $a_0, a_1, \ldots, a_n$ nur auf eine endliche Anzahl von Arten bestimmt werden. Man hat dann unter den so bestimmten Functionen $f(\Theta)$ nur die irreduciblen beizubehalten. Wenn man nun die algebraischen Zahlen in der Weise ordnet, dass man die Zahlen von geringerer Höhe denen von grösserer Höhe voranstellt, dass man unter Zahlen gleicher Höhe die voranstellt, deren reeller Theil kleiner ist, und unter Zahlen von gleicher Höhe und gleichem reellen Theile die von kleinerem imaginären Theile vorangehen lässt, so haben wir eine zählbare Anordnung der algebraischen Zahlen, und es ist erwiesen, dass die Gesammtheit aller algebraischen Zahlen eine abzählbare Menge ist.

Wir erhalten beispielsweise:
\[
\begin{array}{llll}
N=1, & n=1, & a_0=1, & a_1=0\\
N=2, & n=1, & a_0=1, & a_1=\pm 1\\
N=3, & n=1, & a_0=1, & a_1=\pm 2\\
     &      & a_0=2, & a_1=\pm 1\\
     & n=2, & a_0=1, & a_1=0, \quad a_2=1,
\end{array}
\]
und der Anfang der geordneten Reihe der algebraischen Zahlen wird also:
\[
0, -1, +1, -2, -\tfrac{1}{2}, -\sqrt{-1}, \sqrt{-1}, \tfrac{1}{2}, 2, \ldots
\]
